\chapter{Background}
\label{ch:background}

This chapter collects the definitions and mechanisms that the rest of
the thesis takes as given. \Cref{sec:bg-guessing} describes how offline
password guessing works and what a breach exposes.
\Cref{sec:bg-honeywords} states the honeyword primitive introduced by
\textcite{juels2013honeywords} and the HoneyChecker architecture that
supports it. \Cref{sec:bg-flatness} formalizes the classical notion of
flatness relative to a fixed attacker class. \Cref{sec:bg-cnn}
introduces the character-level convolutional networks that
\Cref{ch:related} and \Cref{ch:method} both build on, and
\Cref{sec:bg-fasttext} covers the FastText subword embeddings used by
the HoneyGen family of generators.

Every concept is presented at the level required to follow the
argument in later chapters. The formal set-level flatness definition
that this thesis introduces is deferred to \Cref{sec:method-formalism},
since it is part of the contribution rather than of the background.

\section{Password Guessing After a Breach}
\label{sec:bg-guessing}

Password authentication systems store a transformation of the user's
credential rather than the credential itself. In modern deployments
this transformation is a salted, slow cryptographic hash such as bcrypt,
scrypt, or Argon2. When a user submits a password $w$ at login, the
server computes the hash and compares it to the stored value; a match
authenticates the session.

Two attack models are relevant here. In the \emph{online} model the
attacker submits guesses through the authentication interface, and the
server can rate-limit or lock accounts after a small number of failed
attempts. In the \emph{offline} model the attacker has obtained a copy
of the password file and can test guesses locally at the rate their
hardware allows. Modern GPU rigs push this rate high enough that the
security of the deployment is determined almost entirely by the
password's position in the guessing order used by the attacker.

The guessing order itself has been the subject of a long research
line. \textcite{narayanan2005fast} showed that Markov models of
character sequences accelerate cracking substantially over dictionary
attacks. \textcite{weir2009pcfg} extracted probabilistic context-free
grammars from leaked corpora and produced guesses in decreasing order
of estimated probability. \textcite{melicher2016fast} used neural
language models over character sequences to reach guessing rates that
outperform both Markov and PCFG models on realistic budgets. All three
lines share the same interface: given a candidate password, produce a
score or a rank that reflects how likely it is to appear in a real
distribution.

The consequence for defenders is that a database compromise is silent.
Nothing distinguishes an offline guessing campaign from a period in
which no attacker has the hash file at all, until the attacker's
recovered credentials are used against a live login. Honeywords, the
subject of the next section, are designed to break that silence.

\section{The Honeyword Primitive}
\label{sec:bg-honeywords}

\subsection{Sweetword Storage}

\textcite{juels2013honeywords} propose the following modification to
the password file. For each user account the server stores not a
single hash but a set of $k$ hashes, one of which corresponds to the
real password and the other $k-1$ to \emph{honeywords}: decoy strings
generated to be difficult to distinguish from the real password. We
call the set of $k$ candidates the \emph{sweetword set}, and any
element of it a \emph{sweetword}. The honeywords are not stored in
clear anywhere; only their hashes appear in the password file.

The order of the sweetwords in the file is either randomized per user
or made independent of the identity of the real password. The
information about which position corresponds to the real password is
stored elsewhere, in a component we describe next.

\subsection{The HoneyChecker}

The scheme separates two pieces of information that were previously
combined in a single password file: the credential material (the $k$
hashes) and the identity of the correct one (an index in
$\{1, \ldots, k\}$). The credential material stays on the main
authentication server, and the index is held by a second, hardened
component called the \emph{HoneyChecker}. The HoneyChecker exposes a
narrow interface: given a user identifier and an index, it returns
whether the index is the correct one, and it raises an alarm if the
index is not.

An honest login proceeds as follows. The user submits their password.
The authentication server hashes it and looks for a match among the
$k$ stored sweetword hashes. If a match is found at index $j$, the
server queries the HoneyChecker with that index. The HoneyChecker
returns \emph{accept} when $j$ is the correct index, and the login
succeeds. This is indistinguishable to the user from an ordinary
password check.

A honeyword-triggered login proceeds differently. The attacker, who
has read the password file offline and recovered one of the $k$
hashes, attempts a login using the recovered plaintext. The
authentication server finds a match at some index $j'$ and queries the
HoneyChecker. If $j'$ is not the correct index, the HoneyChecker
returns \emph{reject} and, in addition, raises an alarm. The alarm is
what makes the primitive a breach detector: a decoy login is
observable in real time, whereas a successful compromise otherwise is
not.

The distributed trust assumption is that a single adversary does not
compromise both the password file and the HoneyChecker. Reading only
the password file leaves the attacker with $k$ indistinguishable
hashes; compromising only the HoneyChecker gives away indices with no
credential material to attach them to. This assumption is what makes
the scheme meaningful; \textcite{juels2013honeywords} and the
subsequent survey by \textcite{chakraborty2022hbat} discuss its
practical implications at length, and we take it as given throughout
the thesis.

\subsection{Alarms and Response}

Once the HoneyChecker raises an alarm, the response is a policy
decision outside the scope of this thesis. Depending on the
deployment, the response can lock the affected account, force a
password reset for a class of users, or trigger a wider incident
investigation. The relevant property here is that the alarm is a
signal that the server did not have before the primitive was in place.

\section{Classical Flatness}
\label{sec:bg-flatness}

\subsection{The Attacker's Task}

Given only the sweetword set $S = \{w_1, \ldots, w_k\}$, the attacker
must decide which element is the real password. This is a
$k$-way distinguishing problem. If the honeywords are well
constructed, the attacker's success probability should be close to
$1/k$, the rate of uniform random guessing; if the honeywords are
distinguishable, the success probability is higher.

The exact success probability depends on the attacker's knowledge and
inference capabilities. Juels and Rivest's original paper models the
attacker as a probabilistic ranker: a function $A$ that, given a
candidate password, returns a probability that it comes from the
distribution of real user-chosen passwords. The attacker orders the
$k$ sweetwords by $A$'s output and guesses the top-ranked one.

\subsection{$1/k$-Flatness and $\varepsilon$-Flatness}

Let $A$ be an attacker and let $p_A(S)$ denote the probability that
$A$ identifies $w^{*}$ on its first guess, taken over the random
choice of $S$ and the internal randomness of the generator. The scheme
is called \emph{$1/k$-flat} against $A$ if
%
\begin{equation}
p_A(S) \;=\; 1/k,
\end{equation}
%
which is the uniform-random rate. Since perfect equality is rarely
attainable, the literature relaxes the condition to
\emph{$\varepsilon$-flatness}: the scheme is $\varepsilon$-flat if
%
\begin{equation}
p_A(S) \;\le\; \frac{1 + \varepsilon}{k},
\end{equation}
%
for some fixed $\varepsilon > 0$. The two definitions coincide as
$\varepsilon \to 0$.

Both definitions are attacker-relative in principle: they depend on
the class of function $A$ used. The classical literature restricts $A$
to the probabilistic ranker family described above. \Cref{ch:method}
extends the definition to arbitrary binary discriminators and
formalizes it at the level of the induced distribution over the set;
that extension is one of the contributions of this thesis.

\subsection{Related Security Properties}

\textcite{chakraborty2022hbat} organize the honeyword security
literature around three properties: flatness, denial-of-service (DoS)
resistance, and multiple-system vulnerability (MSV) resistance. DoS
resistance controls the rate at which a legitimate but confused user
can be locked out by accidentally submitting a honeyword. MSV
resistance controls the extent to which learning the sweetwords for
one user helps an attacker for other users. This thesis addresses only
flatness. The DoS and MSV properties are largely orthogonal to the
filtering construction we develop, and we return to them briefly in
\Cref{ch:conclusion}.

\section{Character-Level Convolutional Classifiers}
\label{sec:bg-cnn}

Character-level convolutional networks are the standard architecture
for both attack and defense in this thesis. The design is a small
adaptation of the text classifier introduced by \textcite{kim2014cnn}
and the character-level variants of \textcite{zhang2015character}.

A password $w$ of at most $L$ characters is right-padded to length $L$
and each character is mapped to an embedding vector of dimension $d$
through a lookup table. The resulting $L \times d$ matrix is treated
as a one-dimensional signal with $d$ input channels. One or more
one-dimensional convolutional layers extract features from local
$n$-character windows: each convolution with kernel width $m$ applies
a learned filter to every window of $m$ consecutive positions and
produces a new feature map. A ReLU activation follows each
convolution. A max-pooling step reduces the length dimension, and a
fully connected layer projects the pooled features to the number of
output classes. A final softmax converts the logits to class
probabilities.

For the binary classifier used here, the two output classes are
\emph{real} and \emph{honeyword}, and the quantity of interest is
$P(\text{real} \mid w)$, the probability the network assigns to the
real class for the input password. This score is the raw material for
every flatness measurement in the thesis. The precise architecture and
training procedure are specified in \Cref{sec:method-cnn}.

The property that makes this class of network well suited to password
strings is that its filters operate on local character windows. This
matches the empirical observation that real user-chosen passwords
concentrate distinguishing structure at the level of short substrings
(digit sequences, common suffixes, case patterns), which the
convolutional layers pick up efficiently without requiring token-level
segmentation.

\section{FastText Password Embeddings}
\label{sec:bg-fasttext}

The HoneyGen generator family described in \Cref{sec:method-gens} and
reviewed in \Cref{ch:related} uses FastText embeddings as its
representation of the password space. FastText, introduced by
\textcite{bojanowski2017fasttext}, extends the skipgram word
representation of \textcite{mikolov2013distributed} by representing
each token as the sum of embeddings for its character $n$-grams. A
password like \texttt{qwerty2020} is represented not only by an
embedding for the string as a whole but by the sum of embeddings for
all its character $n$-grams of length in a configured range.

Two properties of this representation matter for honeyword generation.
First, the model produces embeddings for out-of-vocabulary strings,
which is essential because a generator is asked to produce passwords
it has never seen. Second, the nearest neighbors of a password in
embedding space tend to share sub-word structure: a password
ending in \texttt{2020} is close to passwords ending in \texttt{2021},
\texttt{2019}, and so on. This makes the embedding a source of
plausible perturbations of a real password without requiring
hand-written rules.

The HoneyGen family uses a skipgram FastText model trained on
the oracle corpus, with the smallest possible \texttt{minCount} and a
minimum $n$-gram length of two. \Cref{sec:method-gens} gives the exact
configuration used in the experiments; the description here is the
background needed to follow that section.
